\documentclass[11pt,aspectratio=169]{beamer}

\usepackage[utf8]{inputenc}
\usepackage[T1]{fontenc}
\usepackage[french]{babel}
\usepackage{helvet}
\renewcommand{\familydefault}{\sfdefault}
\usepackage{booktabs}
\usepackage{tabularx}
\usepackage{tikz}
\usetikzlibrary{positioning,arrows.meta}

% Palette : bleu #0000FF, encre sombre, fond blanc cassé
\definecolor{bleu}{HTML}{0000FF}
\definecolor{encre}{HTML}{1A1A2E}
\definecolor{gris}{HTML}{5C5C6E}
\definecolor{papier}{HTML}{F7F6F2}

\usetheme{default}
\setbeamertemplate{navigation symbols}{}
\setbeamercolor{background canvas}{bg=papier}
\setbeamercolor{normal text}{fg=encre}
\setbeamercolor{frametitle}{fg=bleu}
\setbeamercolor{title}{fg=bleu}
\setbeamercolor{subtitle}{fg=encre}
\setbeamercolor{structure}{fg=bleu}
\setbeamercolor{alerted text}{fg=bleu}
\setbeamerfont{frametitle}{series=\bfseries,size=\Large}
\setbeamerfont{title}{series=\bfseries,size=\huge}
\setbeamertemplate{itemize items}{\rule[0.25ex]{0.6ex}{0.6ex}}
\setbeamertemplate{itemize subitem}{\textcolor{gris}{\rule[0.3ex]{0.8ex}{0.15ex}}}
\setbeamertemplate{enumerate items}[default]
\setbeamertemplate{frametitle}{%
  \vspace{0.6em}\insertframetitle\par
  \vspace{0.15em}{\color{bleu}\rule{2.2cm}{1.2pt}}}
\setbeamertemplate{footline}{%
  \hspace{1em}\textcolor{gris}{\tiny Littérature sur les portails open data}%
  \hfill\textcolor{gris}{\tiny\insertframenumber/\inserttotalframenumber}\hspace{1em}\vspace{0.6em}}

\newcommand{\refc}[1]{\textcolor{gris}{\footnotesize #1}}

\title{Les portails de données ouvertes}
\subtitle{Ce que dit la littérature, des origines aux portails nationaux africains}
\date{Octobre 2026}

\begin{document}

\begin{frame}[plain]
  \vfill
  {\usebeamerfont{title}\usebeamercolor[fg]{title}\inserttitle\par}
  \vspace{0.6em}
  {\large\insertsubtitle\par}
  \vspace{1.2em}
  {\color{bleu}\rule{4cm}{1.5pt}}\par
  \vspace{0.8em}
  {\small\textcolor{gris}{Revue de littérature pour préparer la construction de portails nationaux}}
  \vfill
\end{frame}

\begin{frame}{Plan}
  \begin{enumerate}
    \item Origines et définitions
    \item Les cadres théoriques : gouvernement ouvert, valeur, écosystèmes
    \item Le portail comme objet de recherche
    \item Qualité, usage et impact
    \item Critiques et limites
    \item Le contexte africain et les pays en développement
    \item Mesurer la maturité : indices et diagnostics
    \item Synthèse et pistes de lecture
  \end{enumerate}
\end{frame}

% ------------------------------------------------------------
\section{Origines}

\begin{frame}{Une généalogie en quatre courants}
  \begin{columns}[T]
    \begin{column}{0.48\textwidth}
      \textbf{\color{bleu}Accès à l'information}
      \begin{itemize}
        \item Lois sur la liberté de l'information (Suède 1766, États-Unis 1966)
        \item Le citoyen a le droit de demander
      \end{itemize}
      \vspace{0.5em}
      \textbf{\color{bleu}Réutilisation des informations publiques}
      \begin{itemize}
        \item Directive européenne PSI de 2003
        \item Logique économique : les données sont une matière première
      \end{itemize}
    \end{column}
    \begin{column}{0.48\textwidth}
      \textbf{\color{bleu}Mouvement du logiciel libre}
      \begin{itemize}
        \item Licences ouvertes, communautés, transparence du code
        \item Open Knowledge Foundation (2004)
      \end{itemize}
      \vspace{0.5em}
      \textbf{\color{bleu}Web de données}
      \begin{itemize}
        \item Berners-Lee : données liées (2006), échelle des 5 étoiles (2010)
        \item Les données doivent être lisibles par les machines
      \end{itemize}
    \end{column}
  \end{columns}
  \vfill
  \refc{Kitchin (2014) ; Davies et al. (2019)}
\end{frame}

\begin{frame}{Les textes fondateurs}
  \small
  \begin{tabularx}{\textwidth}{@{}lX@{}}
    \toprule
    \textbf{Date} & \textbf{Texte et apport} \\
    \midrule
    2007 & \textbf{Principes de Sebastopol} : 8 critères des données publiques ouvertes (complètes, primaires, à jour, accessibles, lisibles par machine, non discriminatoires, non propriétaires, libres de droits) \\
    2009 & \textbf{Mémorandum Obama} et Open Government Directive ; lancement de data.gov \\
    2010 & Lancement de data.gov.uk ; échelle des 5 étoiles \\
    2011 & \textbf{Open Government Partnership} ; Kenya, premier portail national africain \\
    2013 & \textbf{Charte du G8} sur les données ouvertes \\
    2015 & \textbf{Charte internationale des données ouvertes} : ouvertes par défaut \\
    2016 & Principes \textbf{FAIR} pour les données de recherche \\
    2019 & Directive européenne 2019/1024 : données à forte valeur \\
    2022 & Cadre de politique des données de l'Union africaine \\
    \bottomrule
  \end{tabularx}
\end{frame}

\begin{frame}{Définir les données ouvertes}
  \begin{block}{Définition de l'Open Knowledge Foundation}
    Des données que chacun peut librement utiliser, réutiliser et redistribuer, sous réserve tout au plus de citer la source et de partager à l'identique.
  \end{block}
  \vspace{0.5em}
  Trois dimensions reviennent dans toute la littérature :
  \begin{itemize}
    \item \textbf{juridique} : une licence ouverte ;
    \item \textbf{technique} : un format lisible par machine, non propriétaire, téléchargeable en entier ;
    \item \textbf{économique} : gratuité ou coût marginal.
  \end{itemize}
  \vspace{0.5em}
  \textbf{Distinction importante} : \emph{open government data} (données publiques) et \emph{open data} au sens large (recherche, entreprises, société civile).
  \vfill
  \refc{Open Definition 2.1 ; Ubaldi (2013), OCDE}
\end{frame}

% ------------------------------------------------------------
\section{Cadres théoriques}

\begin{frame}{Gouvernement ouvert : trois promesses}
  \begin{columns}[T]
    \begin{column}{0.32\textwidth}
      \textbf{\color{bleu}Transparence}\par\smallskip
      \small Rendre l'action publique visible et contrôlable
    \end{column}
    \begin{column}{0.32\textwidth}
      \textbf{\color{bleu}Participation}\par\smallskip
      \small Associer les citoyens aux décisions
    \end{column}
    \begin{column}{0.32\textwidth}
      \textbf{\color{bleu}Collaboration}\par\smallskip
      \small Co-produire des services avec la société
    \end{column}
  \end{columns}
  \vspace{1em}
  \begin{itemize}
    \item \textbf{Meijer, Curtin et Hillebrandt (2012)} : le gouvernement ouvert combine la \emph{vision} (voir dans l'État) et la \emph{voix} (parler à l'État).
    \item \textbf{Yu et Robinson (2012)} : \og la nouvelle ambiguïté \fg{}. Ouvrir des données techniques (horaires de bus) n'est pas rendre l'État redevable. Il faut distinguer la technologie de la finalité politique.
  \end{itemize}
  \vfill
  \refc{Meijer et al. (2012), IRAS 78(1) ; Yu \& Robinson (2012), UCLA Law Review Discourse 59}
\end{frame}

\begin{frame}{Bénéfices, obstacles et mythes}
  \textbf{Janssen, Charalabidis et Zuiderwijk (2012)}, article le plus cité du champ.
  \vspace{0.4em}
  \begin{columns}[T]
    \begin{column}{0.48\textwidth}
      \textbf{\color{bleu}Bénéfices identifiés}
      \begin{itemize}
        \item politiques et sociaux : transparence, confiance
        \item économiques : innovation, nouveaux services
        \item opérationnels : meilleure décision publique
      \end{itemize}
    \end{column}
    \begin{column}{0.48\textwidth}
      \textbf{\color{bleu}Cinq mythes}
      \begin{enumerate}
        \small
        \item publier produit automatiquement des bénéfices
        \item toute information doit être publiée sans restriction
        \item il suffit de publier
        \item chacun peut utiliser les données
        \item l'open data aboutit à un gouvernement ouvert
      \end{enumerate}
    \end{column}
  \end{columns}
  \vspace{0.5em}
  \alert{Leçon pour un portail} : la publication est une condition nécessaire, jamais suffisante.
  \vfill
  \refc{Janssen et al. (2012), Information Systems Management 29(4)}
\end{frame}

\begin{frame}{La valeur économique}
  \begin{itemize}
    \item \textbf{McKinsey Global Institute (2013)} : 3 000 à 5 000 milliards de dollars de valeur potentielle par an dans 7 secteurs. Chiffre très repris, souvent critiqué pour sa méthode.
    \item \textbf{Rapports de la Commission européenne} (2015, 2020) : estimation de la taille du marché européen des données ouvertes.
    \item \textbf{Ubaldi (2013, OCDE)} : la valeur dépend de la capacité de l'écosystème à réutiliser, pas de la quantité publiée.
  \end{itemize}
  \vspace{0.6em}
  \textbf{Point de consensus} : la valeur est \emph{potentielle} et \emph{indirecte}. Elle se mesure mal, et les estimations macroéconomiques restent fragiles.
  \vfill
  \refc{Manyika et al. (2013) ; Ubaldi (2013), OECD Working Papers on Public Governance 22}
\end{frame}

\begin{frame}{Le tournant des écosystèmes}
  \begin{center}
  \begin{tikzpicture}[node distance=0.6cm and 0.9cm,
      b/.style={draw=bleu,line width=0.8pt,minimum width=2.4cm,minimum height=0.9cm,align=center,font=\small},
      a/.style={-{Stealth},line width=0.8pt,color=encre}]
    \node[b] (prod) {Producteurs\\(administrations)};
    \node[b,right=of prod] (port) {Portail\\(infrastructure)};
    \node[b,right=of port] (inter) {Intermédiaires\\(médias, civic tech)};
    \node[b,right=of inter] (util) {Utilisateurs\\finaux};
    \draw[a] (prod) -- (port);
    \draw[a] (port) -- (inter);
    \draw[a] (inter) -- (util);
    \draw[a,dashed,color=gris] (util.south) to[bend left=20] node[below,font=\scriptsize,color=gris]{retours, demandes de données} (prod.south);
  \end{tikzpicture}
  \end{center}
  \vspace{0.3em}
  \begin{itemize}
    \small
    \item \textbf{Heimstädt, Saunderson et Heath (2014)} : l'écosystème britannique, du \og bambin à l'adolescent \fg{}.
    \item \textbf{Zuiderwijk, Janssen et Davis (2014)} : éléments essentiels d'un écosystème (publication, recherche, usage, retour).
    \item \textbf{Dawes, Vidiasova et Parkhimovich (2016)} : planifier un programme open data comme un écosystème.
  \end{itemize}
  \refc{JeDEM 6(2) ; Information Polity 19 ; Government Information Quarterly 33(1)}
\end{frame}

% ------------------------------------------------------------
\section{Le portail}

\begin{frame}{Le portail comme objet de recherche}
  \begin{itemize}
    \item Première génération (2009 à 2013) : le portail est un \textbf{catalogue} de fichiers.
    \item Deuxième génération : API, visualisations, demandes de données, réutilisations.
    \item Troisième génération : le portail devient une \textbf{plateforme d'interaction} entre producteurs et réutilisateurs.
  \end{itemize}
  \vspace{0.6em}
  Travaux de référence :
  \begin{itemize}
    \small
    \item \textbf{Attard et al. (2015)} : revue systématique des initiatives open data et de leur cycle de vie.
    \item \textbf{Lněnička et Nikiforova (2021)} : \og transparency by design \fg{}, les fonctions qu'un portail doit offrir.
    \item \textbf{Charalabidis et al. (2018)} : \emph{The World of Open Data}, synthèse sur les architectures et les usages.
  \end{itemize}
  \vfill
  \refc{Government Information Quarterly 32(4) ; Telematics and Informatics 61 ; Springer}
\end{frame}

\begin{frame}{Les fonctions attendues d'un portail}
  \small
  \begin{tabularx}{\textwidth}{@{}lX@{}}
    \toprule
    \textbf{Fonction} & \textbf{Éléments relevés dans la littérature} \\
    \midrule
    Découverte & recherche, filtres, catégories, métadonnées riches, moissonnage DCAT \\
    Accès & téléchargement, API, formats ouverts, licences claires \\
    Compréhension & aperçu, visualisation, dictionnaire des données, documentation \\
    Qualité & fraîcheur affichée, schémas, validation, contact du producteur \\
    Interaction & demandes de données, commentaires, signalement d'erreurs \\
    Valorisation & vitrine des réutilisations, statistiques d'usage \\
    \bottomrule
  \end{tabularx}
  \vfill
  \refc{Máchová \& Lněnička (2017) ; Lněnička \& Nikiforova (2021)}
\end{frame}

% ------------------------------------------------------------
\section{Qualité et usage}

\begin{frame}{Évaluer la qualité des portails}
  \begin{itemize}
    \item \textbf{Umbrich, Neumaier et Polleres (2015)} : observatoire automatique de plus de 80 portails CKAN. Constats : métadonnées incomplètes, liens morts, licences absentes.
    \item \textbf{Máchová et Lněnička (2017)} : cadre d'évaluation des portails nationaux sur trois axes, données, fonctions du portail et engagement des utilisateurs.
    \item \textbf{Wilkinson et al. (2016)} : principes FAIR (trouvable, accessible, interopérable, réutilisable), repris au-delà de la recherche.
  \end{itemize}
  \vspace{0.6em}
  \alert{Constat récurrent} : de nombreux portails sont des \og cimetières de données \fg{} : beaucoup de jeux publiés, peu mis à jour, peu utilisés.
  \vfill
  \refc{OBD 2015 ; Journal of Theoretical and Applied Electronic Commerce Research 12(1) ; Scientific Data 3}
\end{frame}

\begin{frame}{Qui utilise les données, et pour quoi ?}
  \begin{itemize}
    \item \textbf{Safarov, Meijer et Grimmelikhuijsen (2017)} : revue systématique de l'utilisation. Types d'usages, conditions, effets, utilisateurs.
    \item \textbf{Ruijer, Grimmelikhuijsen et Meijer (2017)} : l'usage dépend du modèle démocratique visé (contrôle, délibération, participation).
    \item \textbf{Gurstein (2011)} : \og l'\emph{effective use} \fg{}. Sans compétences, accès et outils, l'open data donne du pouvoir à ceux qui en ont déjà.
  \end{itemize}
  \vspace{0.6em}
  \textbf{Implication pour le portail} : concevoir pour des profils précis (journalistes, chercheurs, développeurs, administrations), pas pour un \og public \fg{} abstrait.
  \vfill
  \refc{Information Polity 22(1) ; Government Information Quarterly 34(1) ; First Monday 16(2)}
\end{frame}

\begin{frame}{Politiques nationales comparées}
  \begin{itemize}
    \item \textbf{Huijboom et Van den Broek (2011)} : première comparaison internationale des stratégies (États-Unis, Royaume-Uni, Australie, Danemark, Espagne).
    \item \textbf{Zuiderwijk et Janssen (2014)} : grille de comparaison des politiques open data (contexte, politique, conditions, impacts).
    \item \textbf{Ubaldi (2013, OCDE)} : cadre d'analyse pour les gouvernements, repris dans les enquêtes OURdata de l'OCDE.
  \end{itemize}
  \vspace{0.6em}
  \textbf{Facteurs de succès les plus cités} : portage politique, cadre juridique, gouvernance des données, financement durable, implication de la demande.
  \vfill
  \refc{European Journal of ePractice 12 ; Government Information Quarterly 31(1)}
\end{frame}

% ------------------------------------------------------------
\section{Critiques}

\begin{frame}{Les critiques à connaître}
  \begin{itemize}
    \item \textbf{Gurstein (2011)} : risque d'accroître les inégalités, au profit des acteurs déjà équipés.
    \item \textbf{Johnson (2014)} : de l'open data à la \og justice informationnelle \fg{}. Les données encodent des choix et des biais de ceux qui les produisent.
    \item \textbf{Yu et Robinson (2012)} : confusion entre ouverture technique et redevabilité politique.
    \item \textbf{Kitchin (2014)} : marchandisation possible des données publiques par de grands acteurs privés.
    \item \textbf{Vie privée} : risques de réidentification lors du croisement de jeux de données anonymisés.
  \end{itemize}
  \vspace{0.5em}
  \alert{Question à se poser pour chaque portail} : à qui profite l'ouverture, et qui en est exclu ?
  \vfill
  \refc{Ethics and Information Technology 16(4) ; Kitchin, \emph{The Data Revolution}, Sage}
\end{frame}

% ------------------------------------------------------------
\section{Afrique}

\begin{frame}{L'open data dans les pays en développement}
  \begin{itemize}
    \item \textbf{Davies et al. (2019)}, \emph{The State of Open Data} : bilan mondial en 30 chapitres, publié par African Minds et le CRDI. Lecture prioritaire.
    \item \textbf{Open Data in Developing Economies} (réseau ODDC, 2013 à 2015) : études de cas montrant que le contexte institutionnel compte plus que la technologie.
    \item \textbf{Constats} : portails lancés avec l'appui de bailleurs puis abandonnés ; faible demande locale ; problèmes de connectivité ; manque de données numérisées en amont.
  \end{itemize}
  \vfill
  \refc{Davies, Walker, Rubinstein \& Perini (dir.), 2019, African Minds \& IDRC}
\end{frame}

\begin{frame}{Travaux sur l'Afrique}
  \begin{itemize}
    \item \textbf{Mutuku et Colaco (2012)} : le portail du Kenya, premier d'Afrique subsaharienne (2011). Forte publication, faible consommation locale.
    \item \textbf{van Schalkwyk, Willmers et McNaughton (2016)} : les données \og visqueuses \fg{}. En Afrique du Sud, les \textbf{intermédiaires} (universités, médias, ONG) sont indispensables pour que les données circulent.
    \item \textbf{Banque mondiale, rapports ODRA} : diagnostics nationaux (plusieurs pays africains), qui pointent surtout la gouvernance et le financement.
    \item \textbf{Banque africaine de développement, Open Data for Africa} : réseau de portails statistiques nationaux, centré sur les indicateurs et les instituts de statistique.
    \item \textbf{Union africaine (2022)} : cadre continental de politique des données, articulé avec la Convention de Malabo.
  \end{itemize}
  \vfill
  \refc{ICEGOV 2012 ; Information Technology for Development 22(sup1)}
\end{frame}

\begin{frame}{Spécificités africaines dégagées par la littérature}
  \small
  \begin{tabularx}{\textwidth}{@{}lX@{}}
    \toprule
    \textbf{Enjeu} & \textbf{Ce que montrent les travaux} \\
    \midrule
    Offre & les données existent souvent sur papier ou dans des fichiers dispersés \\
    Demande & faible, mais forte chez les journalistes, ONG et chercheurs quand on les implique \\
    Intermédiaires & rôle central de Code for Africa, des universités et des médias \\
    Statistique & l'institut national de la statistique est l'acteur clé, d'où des portails orientés indicateurs \\
    Financement & dépendance aux bailleurs, risque d'arrêt après le projet \\
    Infrastructure & connectivité, coût des données mobiles, hébergement \\
    \bottomrule
  \end{tabularx}
\end{frame}

% ------------------------------------------------------------
\section{Mesure}

\begin{frame}{Mesurer la maturité}
  \small
  \begin{tabularx}{\textwidth}{@{}lX@{}}
    \toprule
    \textbf{Outil} & \textbf{Ce qu'il mesure} \\
    \midrule
    ODRA (Banque mondiale) & préparation d'un pays, en 8 dimensions, avant le lancement \\
    Open Data Barometer (Web Foundation, 2013 à 2017) & préparation, mise en œuvre et impact \\
    Global Data Barometer (2022) & gouvernance, capacités, disponibilité et usage des données \\
    ODIN (Open Data Watch) & ouverture et couverture des statistiques officielles \\
    OURdata (OCDE) & disponibilité, accessibilité, soutien à la réutilisation \\
    Open Data Maturity (UE) & politique, portail, qualité et impact en Europe \\
    \bottomrule
  \end{tabularx}
  \vspace{0.6em}
  Ces indices sont aussi des \textbf{instruments de politique} : ils orientent ce que les pays choisissent de publier.
\end{frame}

% ------------------------------------------------------------
\section{Synthèse}

\begin{frame}{Ce qu'il faut retenir}
  \begin{enumerate}
    \item Un portail est une \textbf{infrastructure au service d'un écosystème}, pas une fin en soi.
    \item La valeur naît de la \textbf{réutilisation}, qui dépend des compétences et des intermédiaires.
    \item La \textbf{qualité et la fraîcheur} comptent plus que le nombre de jeux publiés.
    \item La réussite dépend d'abord de la \textbf{gouvernance}, du \textbf{cadre juridique} et du \textbf{financement durable}.
    \item En Afrique, l'institut de statistique et les intermédiaires sont au centre du dispositif.
    \item Il faut garder un regard critique : \textbf{qui gagne}, qui est exclu, quels risques pour la vie privée.
  \end{enumerate}
\end{frame}

\begin{frame}{Ordre de lecture conseillé}
  \begin{enumerate}
    \small
    \item Janssen, Charalabidis et Zuiderwijk (2012) : les bénéfices et les mythes
    \item Ubaldi (2013, OCDE) : le cadre d'analyse pour les gouvernements
    \item Gurstein (2011) puis Johnson (2014) : les critiques
    \item Attard et al. (2015) : la revue systématique
    \item Zuiderwijk, Janssen et Davis (2014) puis Dawes et al. (2016) : les écosystèmes
    \item Máchová et Lněnička (2017) : évaluer un portail
    \item Davies et al. (2019) : \emph{The State of Open Data}, en entier ou par chapitres
    \item van Schalkwyk et al. (2016) : les intermédiaires en Afrique
    \item Banque mondiale, Open Data Toolkit et un rapport ODRA africain
  \end{enumerate}
\end{frame}

\begin{frame}[allowframebreaks]{Références}
  \scriptsize
  \begin{itemize}
    \item Attard, J., Orlandi, F., Scerri, S., Auer, S. (2015). A systematic review of open government data initiatives. \emph{Government Information Quarterly}, 32(4).
    \item Charalabidis, Y., Zuiderwijk, A., Alexopoulos, C., Janssen, M., Lampoltshammer, T., Ferro, E. (2018). \emph{The World of Open Data}. Springer.
    \item Davies, T., Walker, S., Rubinstein, M., Perini, F. (dir.) (2019). \emph{The State of Open Data: Histories and Horizons}. African Minds et IDRC.
    \item Dawes, S. S., Vidiasova, L., Parkhimovich, O. (2016). Planning and designing open government data programs: An ecosystem approach. \emph{Government Information Quarterly}, 33(1).
    \item Gurstein, M. (2011). Open data: Empowering the empowered or effective data use for everyone? \emph{First Monday}, 16(2).
    \item Heimstädt, M., Saunderson, F., Heath, T. (2014). From toddler to teen: Growth of an open data ecosystem. \emph{JeDEM}, 6(2).
    \item Huijboom, N., Van den Broek, T. (2011). Open data: an international comparison of strategies. \emph{European Journal of ePractice}, 12.
    \item Janssen, M., Charalabidis, Y., Zuiderwijk, A. (2012). Benefits, adoption barriers and myths of open data and open government. \emph{Information Systems Management}, 29(4).
    \item Johnson, J. A. (2014). From open data to information justice. \emph{Ethics and Information Technology}, 16(4).
    \item Kitchin, R. (2014). \emph{The Data Revolution}. Sage.
    \item Lněnička, M., Nikiforova, A. (2021). Transparency-by-design: What is the role of open data portals? \emph{Telematics and Informatics}, 61.
    \item Máchová, R., Lněnička, M. (2017). Evaluating the quality of open data portals on the national level. \emph{Journal of Theoretical and Applied Electronic Commerce Research}, 12(1).
    \item Manyika, J. et al. (2013). \emph{Open data: Unlocking innovation and performance with liquid information}. McKinsey Global Institute.
    \item Meijer, A., Curtin, D., Hillebrandt, M. (2012). Open government: connecting vision and voice. \emph{International Review of Administrative Sciences}, 78(1).
    \item Mutuku, L., Colaco, J. (2012). Increasing Kenyan open data consumption. \emph{ICEGOV 2012}.
    \item Ruijer, E., Grimmelikhuijsen, S., Meijer, A. (2017). Open data for democracy. \emph{Government Information Quarterly}, 34(1).
    \item Safarov, I., Meijer, A., Grimmelikhuijsen, S. (2017). Utilization of open government data: A systematic literature review. \emph{Information Polity}, 22(1).
    \item Ubaldi, B. (2013). Open government data: Towards empirical analysis of open government data initiatives. \emph{OECD Working Papers on Public Governance}, 22.
    \item Umbrich, J., Neumaier, S., Polleres, A. (2015). Quality assessment and evolution of open data portals. \emph{OBD 2015}.
    \item van Schalkwyk, F., Willmers, M., McNaughton, M. (2016). Viscous open data: The roles of intermediaries in an open data ecosystem. \emph{Information Technology for Development}, 22(sup1).
    \item Wilkinson, M. D. et al. (2016). The FAIR Guiding Principles for scientific data management and stewardship. \emph{Scientific Data}, 3.
    \item Yu, H., Robinson, D. G. (2012). The new ambiguity of open government. \emph{UCLA Law Review Discourse}, 59.
    \item Zuiderwijk, A., Janssen, M. (2014). Open data policies, their implementation and impact. \emph{Government Information Quarterly}, 31(1).
    \item Zuiderwijk, A., Janssen, M., Davis, C. (2014). Innovation with open data: Essential elements of open data ecosystems. \emph{Information Polity}, 19.
  \end{itemize}
\end{frame}

\end{document}
